\documentclass[11pt,a4paper]{article}
\usepackage[utf8]{inputenc}
\usepackage[T1]{fontenc}
\usepackage[portuguese]{babel}
\usepackage{geometry}
\geometry{margin=2.5cm}
\usepackage{listings}
\usepackage{xcolor}
\usepackage{amsmath}
\usepackage{amssymb}
\usepackage{hyperref}
\usepackage{enumitem}
\usepackage{tcolorbox}
\usepackage{tabularx}
\usepackage{booktabs}

\definecolor{codebg}{rgb}{0.95,0.95,0.95}
\definecolor{codegreen}{rgb}{0.1,0.5,0.1}
\definecolor{codepurple}{rgb}{0.5,0.1,0.5}
\definecolor{codegray}{rgb}{0.5,0.5,0.5}

\lstdefinestyle{python}{
    backgroundcolor=\color{codebg},
    commentstyle=\color{codegray}\itshape,
    keywordstyle=\color{codepurple}\bfseries,
    stringstyle=\color{codegreen},
    basicstyle=\ttfamily\small,
    breaklines=true,
    frame=single,
    language=Python,
    showstringspaces=false,
    tabsize=4,
    literate={á}{{\'a}}1 {ã}{{\~a}}1 {é}{{\'e}}1 {í}{{\'i}}1 {ó}{{\'o}}1 {õ}{{\~o}}1 {ú}{{\'u}}1 {ç}{{\c{c}}}1 {Á}{{\'A}}1 {Ã}{{\~A}}1 {É}{{\'E}}1 {Í}{{\'I}}1 {Ó}{{\'O}}1 {Õ}{{\~O}}1 {Ú}{{\'U}}1 {Ç}{{\c{C}}}1 {â}{{\^a}}1 {ê}{{\^e}}1 {ô}{{\^o}}1 {û}{{\^u}}1 {Â}{{\^A}}1 {Ê}{{\^E}}1 {Ô}{{\^O}}1 {Û}{{\^U}}1 {à}{{\`a}}1 {è}{{\`e}}1 {ì}{{\`i}}1 {ò}{{\`o}}1 {ù}{{\`u}}1
}
\lstset{style=python}

\newtcolorbox{objetivos}[1][]{
  colback=green!5!white,
  colframe=green!40!black,
  title=O que você vai aprender nesta página,
  fonttitle=\bfseries,
  #1
}

\title{\textbf{Data Analysis with Python 2024--2025}\\Parte 1: Python\\\large Conteúdo Teórico}
\author{Tradução e Adaptação: Universidade de Helsinque (MOOC.fi)}
\date{}

\begin{document}
\maketitle

\section*{Nota Importante}
Este material é uma tradução e adaptação baseada no curso \textit{Data Analysis with Python 2024-2025}, oferecido pela \textbf{Universidade de Helsinque (MOOC.fi)}.

\begin{objetivos}
\begin{itemize}
    \item Aprender os conceitos básicos da linguagem Python.
    \item Compreender strings e módulos.
    \item Trabalhar com sequências e objetos iteráveis.
\end{itemize}
\end{objetivos}

\tableofcontents
\newpage

\section{Entrada e saída básicas}
O programa clássico ``Hello, world!'' é extremamente simples em Python. A função \texttt{print()} envia texto para a saída padrão. Vários argumentos podem ser passados para \texttt{print}; por padrão, eles são separados por espaços. Expressões numéricas são avaliadas antes de serem impressas.

\begin{lstlisting}
print("Hello world2!")
print("Hello,", "John!", "How are you?")
print(1, "plus", 2, "equals", 1+2)
\end{lstlisting}

A função \texttt{input()} permite ler texto digitado pelo usuário. O resultado de \texttt{input()} é uma string.
\begin{lstlisting}
name = input("Give me your name: ")
print("Hello,", name)
\end{lstlisting}

\subsection{Indentação}
Python usa indentação para marcar blocos de código. Quatro espaços são uma convenção comum. Diferentemente de linguagens que usam chaves, o fim da indentação marca o fim do bloco.

\begin{lstlisting}
for i in range(3):
    print("Hello")
print("Bye!")
\end{lstlisting}

A expressão \texttt{range(n)} produz os inteiros de 0 a \(n-1\); o limite final não é incluído.

\section{Variáveis e tipos de dados}
Uma atribuição associa um nome a um valor. Python não exige a declaração prévia de uma variável nem a especificação do tipo. O tipo pertence ao objeto referenciado, e isso caracteriza a tipagem dinâmica.

Os tipos básicos destacados neste material incluem \texttt{int}, \texttt{float}, \texttt{complex}, \texttt{str}, \texttt{bool} e \texttt{bytes}. A função \texttt{type()} permite descobrir o tipo do valor atual de um nome.

\begin{lstlisting}
a = 1
print(type(a))
a = "algum texto"
print(type(a))
\end{lstlisting}

Os construtores de tipo também podem ser usados para conversão, por exemplo \texttt{int()}, \texttt{float()}, \texttt{str()} e \texttt{bool()}.

Bytes representam unidades de informação de 8 bits. Em situações que envolvem arquivos ou conjuntos de dados, a codificação de texto pode precisar ser especificada. O exemplo de UTF-8 mostra que um caractere pode ocupar mais de um byte.

\section{Strings}
Strings são sequências de caracteres delimitadas por aspas simples ou duplas. Aspas podem ser escapadas com barra invertida, e sequências como \texttt{\textbackslash n} e \texttt{\textbackslash t} representam nova linha e tabulação.

Strings também podem ser escritas em múltiplas linhas usando aspas triplas. Embora seja possível concatenar com \texttt{+}, para muitas strings o método \texttt{join()} é preferível por eficiência.

\subsection{Interpolação de strings}
O material apresenta três formas: formatação tradicional, o método \texttt{format()} e f-strings. Atualmente, f-strings e \texttt{format()} são as formas mais usadas nos exemplos.

\begin{lstlisting}
print("{} plus {} is equal to {}".format(1, 3, 4))
print(f"{1} plus {3} is equal to {4}")
print(f"{1.6:.1f} {1.7:.2f} {1.8:.3f}")
\end{lstlisting}

Especificadores permitem controlar casas decimais e largura de campos. Para strings, o especificador \texttt{s} pode ser usado; para inteiros, \texttt{d}; e para números de ponto flutuante, \texttt{f}.

\section{Expressões e instruções}
Uma expressão é um trecho de código que produz um valor. Exemplos incluem operações aritméticas, chamadas de funções, indexação, comparações e acesso a atributos.

Uma instrução é um comando que produz um efeito. Atribuições e chamadas de função isoladas são exemplos. Python possui operadores de atribuição aumentada como \texttt{+=}, \texttt{-=}, \texttt{*=}, \texttt{/=} e outros.

\section{Laços para tarefas repetitivas}
Python possui principalmente os laços \texttt{while} e \texttt{for}. Um \texttt{while} repete enquanto uma condição for verdadeira.

\begin{lstlisting}
i = 1
while i*i < 1000:
    print("Square of", i, "is", i*i)
    i += 1
\end{lstlisting}

O \texttt{for} percorre os elementos de um iterável. Listas e objetos criados por \texttt{range()} são exemplos.

\begin{lstlisting}
s = 0
for i in [0,1,2,3,4,5,6,7,8,9]:
    s += i
print("The sum is", s)
\end{lstlisting}

\section{Decisões com if}
A estrutura \texttt{if/elif/else} seleciona blocos de acordo com condições booleanas.

\begin{lstlisting}
x = int(input("Give an integer: "))
if x >= 0:
    a = x
else:
    a = -x
print("The absolute value of %i is %i" % (x, a))
\end{lstlisting}

\subsection{break e continue}
\texttt{break} encerra o laço atual quando a condição desejada é encontrada. \texttt{continue} interrompe apenas a iteração corrente e passa para a próxima.

\section{Funções}
Funções são declaradas com \texttt{def}. Uma função pode possuir parâmetros, retornar um valor e conter uma docstring descrevendo sua finalidade.

\begin{lstlisting}
def double(x):
    # Multiplica o argumento por dois.
    return x*2

print(double(4))
\end{lstlisting}

Python permite parâmetros posicionais, parâmetros nomeados, valores padrão e empacotamento/desempacotamento com \texttt{*} e \texttt{**}. Funções criam um novo escopo local.

\subsection{Escopo}
Variáveis definidas dentro de uma função são locais. Variáveis globais são acessíveis de dentro da função, mas uma atribuição normalmente cria um nome local. A instrução \texttt{global} pode ser usada quando é realmente necessário reatribuir um global; em funções aninhadas, \texttt{nonlocal} referencia o escopo externo mais próximo.

\section{Estruturas de dados}
As principais estruturas destacadas são strings, listas, tuplas, dicionários e conjuntos.

\subsection{Sequências}
Listas podem ter quantidade arbitrária de elementos e são mutáveis. Tuplas são sequências ordenadas e imutáveis. Strings, listas e tuplas compartilham operações como \texttt{len()}, indexação, fatiamento, concatenação e repetição.

A indexação começa em zero. Índices negativos contam a partir do final. Slices seguem a forma \texttt{sequencia[inicio:fim:passo]}, com o limite final excluído.

\subsection{Modificação de listas}
Listas podem ser alteradas por indexação, slicing e métodos como \texttt{append}, \texttt{extend}, \texttt{insert}, \texttt{remove}, \texttt{pop}, \texttt{reverse} e \texttt{sort}. A função \texttt{sorted()} devolve uma nova lista ordenada sem modificar a original.

\subsection{range e ordenação}
\texttt{range()} produz uma sequência numérica eficiente sem materializar obrigatoriamente uma lista. Seu passo pode ser especificado. Para ordenar, \texttt{sort()} altera a lista e \texttt{sorted()} produz outra.

\subsection{zip}
\texttt{zip()} combina várias sequências. Se duas sequências forem combinadas, os elementos resultam em pares; com \(n\) sequências, o resultado tem tuplas de \(n\) elementos. O tamanho é limitado pela sequência mais curta.

\subsection{enumerate}
\texttt{enumerate()} fornece simultaneamente índice e elemento. Conceitualmente, é equivalente a combinar \texttt{range(len(L))} e \texttt{L} com \texttt{zip()}.

\subsection{Dicionários}
Dicionários associam chaves a valores. As chaves precisam ser hashable. A consulta é feita por indexação. Entre os métodos importantes estão \texttt{items()}, \texttt{keys()}, \texttt{values()}, \texttt{get()}, \texttt{update()}, \texttt{pop()} e \texttt{popitem()}.

\subsection{Conjuntos}
Conjuntos armazenam elementos únicos e suportam união, interseção, diferença e diferença simétrica. Os operadores correspondentes são \texttt{|}, \texttt{\&}, \texttt{-} e \texttt{\string^}.

\subsection{Operador in, desempacotamento e del}
O operador \texttt{in} testa pertinência em contêineres. Sequências e dicionários também podem ser desempacotados em múltiplos nomes. A instrução \texttt{del} remove uma associação ou um item/slice de um contêiner.

\section{Comprehensions}
Comprehensions permitem construir estruturas de dados de maneira compacta. O material apresenta list comprehensions, dictionary comprehensions e set comprehensions, incluindo múltiplos laços e filtros.

\section{Processamento de sequências}
Funções podem ser tratadas como objetos de primeira classe: podem ser passadas como argumentos, retornadas por funções e armazenadas em variáveis ou estruturas de dados. Nesta parte são estudadas \texttt{map}, \texttt{filter}, \texttt{reduce} e expressões \texttt{lambda}.

\subsection{map e lambda}
\texttt{map()} aplica uma função a cada elemento e devolve um iterável. Em Python moderno, é comum convertê-lo explicitamente em lista quando se deseja visualizar ou armazenar os resultados.

\begin{lstlisting}
def double(x):
    return 2*x
L = [12, 4, -1]
print(list(map(double, L)))
\end{lstlisting}

Uma \texttt{lambda} define uma pequena função sem nome.
\begin{lstlisting}
L = [2, 3, 5]
print(list(map(lambda x: 2*x+x**2, L)))
\end{lstlisting}

\subsection{filter}
\texttt{filter()} conserva apenas os elementos para os quais a função predicado retorna \texttt{True}. O material observa que comprehensions são frequentemente mais flexíveis em Python moderno.

\subsection{reduce}
\texttt{reduce()} combina iterativamente elementos de um iterável até obter um único resultado. Em Python, está em \texttt{functools}. Pode receber um valor inicial.

\section{Manipulação de strings}
Strings são imutáveis. Seus métodos não alteram a string original; em vez disso, retornam novos valores. O material organiza os métodos em grupos: classificação, transformação, busca, ajuste e junção/divisão.

\subsection{Classificação}
Métodos como \texttt{isalpha()}, \texttt{isdigit()}, \texttt{isalnum()}, \texttt{islower()}, \texttt{isupper()} e \texttt{isspace()} produzem valores booleanos.

\subsection{Transformação}
\texttt{lower()}, \texttt{upper()}, \texttt{capitalize()}, \texttt{title()} e \texttt{swapcase()} são exemplos de transformação de texto.

\subsection{Busca}
Métodos como \texttt{find()}, \texttt{index()}, \texttt{count()}, \texttt{startswith()} e \texttt{endswith()} ajudam a localizar substrings.

\subsection{Ajuste}
\texttt{strip()}, \texttt{lstrip()} e \texttt{rstrip()} removem caracteres das extremidades. O argumento pode especificar um conjunto de caracteres.

\subsection{join e split}
\texttt{join()} combina uma sequência de strings usando um separador. Para muitas strings, ele é mais eficiente que repetidas concatenações com \texttt{+}. \texttt{split()} divide uma string em partes, por um separador ou por espaços em branco quando nenhum separador é fornecido.

\section{Módulos}
Módulos dividem programas grandes em unidades menores e independentes. Cada arquivo \texttt{.py} pode funcionar como um módulo. A biblioteca padrão possui centenas de módulos; entre os exemplos citados estão \texttt{re}, \texttt{math}, \texttt{random}, \texttt{os} e \texttt{sys}.

\subsection{Usando módulos}
A forma \texttt{import math} permite acessar funções com o nome do módulo, como \texttt{math.cos(0)}. Também é possível importar nomes específicos com \texttt{from math import cos}. Embora \texttt{from math import *} exista, essa forma pode causar confusão por importar muitos nomes para o escopo atual.

\subsection{Localização de módulos}
Ao importar um módulo, Python procura primeiro módulos embutidos e depois os caminhos de \texttt{sys.path}. Isso explica por que arquivos locais podem ser importados e por que colisões de nomes podem ocorrer.

\subsection{Hierarquia de módulos e pacotes}
Pacotes são módulos que podem conter outros módulos e subpacotes. A hierarquia do sistema de arquivos representa a hierarquia de importação. Historicamente, arquivos \texttt{\_\_init\_\_.py} são usados para compor um pacote.

\subsection{Conteúdo de um módulo}
Atribuições, classes e funções definidas em um arquivo tornam-se atributos do módulo. O atributo \texttt{\_\_name\_\_} vale \texttt{"\_\_main\_\_"} quando o arquivo é executado diretamente. Isso permite usar o padrão:

\begin{lstlisting}
def main():
    for _ in range(3):
        print("Hello")

if __name__ == "__main__":
    main()
\end{lstlisting}

\section{Resumo}
Os blocos de código em Python são definidos por indentação consistente. O laço \texttt{for} percorre os elementos de um iterável sem exigir o controle explícito dos índices. Python possui tipagem dinâmica, bom suporte para manipulação de strings e uma grande variedade de ferramentas para processar sequências, como \texttt{map}, \texttt{filter}, \texttt{reduce}, \texttt{zip}, \texttt{enumerate} e \texttt{range}. O material também destaca o uso de \texttt{join()} para concatenar muitas strings e o papel das funções anônimas \texttt{lambda}.
\end{document}
